\documentclass[11pt]{article}
\usepackage{amsmath,amssymb,amsthm,mathtools}
\usepackage[margin=1in]{geometry}
\usepackage{booktabs}
\usepackage{enumitem}

\newtheorem{theorem}{Theorem}
\newtheorem{lemma}{Lemma}
\newtheorem{proposition}{Proposition}
\newtheorem{definition}{Definition}
\newtheorem{warning}{Warning}
\newtheorem{remark}{Remark}

\newcommand{\R}{\mathbb R}
\newcommand{\C}{\mathbb C}
\newcommand{\ip}[2]{\langle #1,#2\rangle}
\newcommand{\norm}[1]{\left\lVert #1\right\rVert}
\newcommand{\pre}{\preceq}

\title{Step 80: Gamma/Archimedean Feature Matching for the RH Membrane Route}
\author{Working note}
\date{}

\begin{document}
\maketitle

\section{Purpose}
Step 79 isolated the completion-balance gate: the signed prime term becomes a positive shift feature only after a diagonal repair. The next analytic question is whether the archimedean/gamma term of the completed zeta package admits a positive feature-map representation compatible with that repair. This note records the operator-theoretic gate that must be earned.

The result here is not a proof of RH and not a proof of the full Douglas bridge. It shows that the gamma slot has a natural positive difference-form representation after normalization, and it states precisely which constants, null modes, and completion terms remain to be matched.

\section{Digamma difference as a positive difference symbol}

\begin{lemma}[Digamma difference formula]
Let $a>0$. For $x\in\R$,
\begin{equation}
\Re\psi(a+ix)-\psi(a)
=\int_0^\infty {e^{-at}\over 1-e^{-t}}\bigl(1-\cos(xt)\bigr)\,dt,
\end{equation}
where $\psi=\Gamma'/\Gamma$ is the digamma function.
Consequently
\begin{equation}
\Psi_a(x):=\Re\psi(a+ix)-\psi(a)\ge0,
\qquad
\Psi_a(0)=0.
\end{equation}
\end{lemma}

\begin{proof}
Use the standard integral representation, valid for $\Re z>0$,
\begin{equation}
\psi(z)=-\gamma+\int_0^\infty\left({e^{-t}\over t}-{e^{-zt}\over 1-e^{-t}}\right)dt.
\end{equation}
Subtract the expression at $z=a$, take real parts, and obtain the stated identity.
\end{proof}

\begin{definition}[Archimedean zeta symbol]
For the completed zeta factor $\pi^{-s/2}\Gamma(s/2)$ at the centered critical-line spectral coordinate, define the normalized archimedean symbol
\begin{equation}
\Psi_\infty(\xi)
:=\Re\psi\left({1\over4}+{i\xi\over2}\right)-\psi\left({1\over4}\right).
\end{equation}
Equivalently,
\begin{equation}
\Psi_\infty(\xi)
=\int_0^\infty {e^{-t/4}\over 1-e^{-t}}\left(1-\cos\left({\xi t\over2}\right)\right)dt.
\end{equation}
\end{definition}

\begin{proposition}[Positive archimedean feature]
The form
\begin{equation}
q_\infty^+[f]
=\int_\R \Psi_\infty(\xi)|\widehat f(\xi)|^2\,d\xi
\end{equation}
is positive on every Fourier-side domain for which the integral is finite. Equivalently, it is generated by a positive shift-difference feature with kernel
\begin{equation}
\kappa_\infty(t)={1\over 2}{e^{-t/4}\over 1-e^{-t}},\qquad t>0,
\end{equation}
with shifts at scale $t/2$:
\begin{equation}
q_\infty^+[f]
=\int_0^\infty \kappa_\infty(t)\,\norm{\tau_{t/2}f-f}_{2}^{2}\,dt.
\end{equation}
\end{proposition}

\begin{proof}
By Plancherel,
\begin{equation}
\norm{\tau_{t/2}f-f}_{2}^{2}
=\int_\R |e^{i\xi t/2}-1|^2|\widehat f(\xi)|^2\,d\xi
=2\int_\R(1-\cos(\xi t/2))|\widehat f(\xi)|^2\,d\xi.
\end{equation}
Multiplying by $\kappa_\infty(t)$ and integrating in $t$ gives the claimed symbol.
\end{proof}

\begin{warning}[Renormalized paired form, not free diagonal]
Near $t=0$,
\begin{equation}
\kappa_\infty(t)\sim {1\over 2t}.
\end{equation}
Thus the separated diagonal $\int\kappa_\infty(t)dt$ diverges near zero. The paired difference form can still be finite on smooth cores because $\norm{\tau_{t/2}f-f}_2^2=O(t^2)$ as $t\to0$. Therefore the archimedean slot may contribute a positive renormalized difference form, but its divergent separated diagonal cannot be spent as an independent finite diagonal payment.
\end{warning}

\section{Gamma matching gate}
The signed archimedean term appearing in a chosen form of the explicit formula must be matched to the positive symbol above plus constants, pole terms, and null-mode records. Abstractly, write
\begin{equation}
R_{\infty,L}^{\rm sgn}[f]
= q_\infty^+[f]-c_{\infty,L}\norm f^2 + r_{\infty,L}[f],
\end{equation}
where $r_{\infty,L}$ is the residual caused by normalization, support restrictions, or a mismatch between the chosen explicit formula convention and the positive gamma feature. The gamma matching gate is:
\begin{equation}
\boxed{\quad r_{\infty,L}\text{ is either zero, positive, or carried by an explicit defect }E_{\infty,L}.\quad}
\end{equation}
If $r_{\infty,L}$ has an unbudgeted negative part, the completed Douglas bridge is not accepted.

\section{Completion balance after gamma matching}
From the prime slot one has
\begin{equation}
q_{{\rm pr},L}[f]
=P_{{\rm pr},L}(h_f)+d_{{\rm pr},L}\norm f^2,
\qquad
 d_{{\rm pr},L}=2\sum_a w_{a,L}.
\end{equation}
After inserting the positive gamma feature, the completion-balance form becomes
\begin{align}
\mathcal B_L[f]
={}&d_{{\rm pr},L}\norm f^2
+q_\infty^+[f]
+q_{{\rm pole},L}^+[f]
+q_{{\rm tail},L}^+[f]
-R_{\infty,L}^{\rm sgn}[f]
-R_{{\rm pole},L}^{\rm sgn}[f]
-R_{{\rm tail},L}^{\rm sgn}[f].
\end{align}
If $R_{\infty,L}^{\rm sgn}=q_\infty^+-c_{\infty,L}\norm f^2+r_{\infty,L}$, then
\begin{align}
\mathcal B_L[f]
={}&(d_{{\rm pr},L}+c_{\infty,L})\norm f^2
+q_{{\rm pole},L}^+[f]
+q_{{\rm tail},L}^+[f]
-r_{\infty,L}[f]
-R_{{\rm pole},L}^{\rm sgn}[f]
-R_{{\rm tail},L}^{\rm sgn}[f].
\end{align}
Thus the gamma feature cancels its paired positive component; what remains is a finite bookkeeping problem involving constants, pole/completion, residuals, and tails.

\begin{theorem}[Gamma-slot positive matching gate]
Suppose on a dense anti-invariant log-side core $\mathcal C^-_{\log}$:
\begin{enumerate}[label=(\roman*)]
\item the signed archimedean term admits the decomposition
$R_{\infty,L}^{\rm sgn}=q_\infty^+-c_{\infty,L}\norm\cdot^2+r_{\infty,L}$;
\item the residual has defect bound $r_{\infty,L}^{-}\preceq E_{\infty,L}$;
\item the remaining pole/tail completion balance satisfies
\begin{equation}
(d_{{\rm pr},L}+c_{\infty,L})\norm\cdot^2+q_{{\rm pole},L}^++q_{{\rm tail},L}^+-R_{{\rm pole},L}^{\rm sgn}-R_{{\rm tail},L}^{\rm sgn}\succeq -E_{{\rm pole/tail},L}.
\end{equation}
\end{enumerate}
Then
\begin{equation}
\mathcal B_L\succeq -(E_{\infty,L}+E_{{\rm pole/tail},L}).
\end{equation}
Consequently any signed Weil inequality
\begin{equation}
q_Z\le q_{{\rm sgn},L}+E_{{\rm Weil},L}
\end{equation}
implies the positive-feature inequality
\begin{equation}
q_Z\le q_{{\rm cmp},L}+E_{{\rm Weil},L}+E_{\infty,L}+E_{{\rm pole/tail},L}.
\end{equation}
If the total defect vanishes and the core-to-full gate passes, then
\begin{equation}
V_Z^*V_Z\preceq W_{{\rm cmp},L}^*W_{{\rm cmp},L},
\end{equation}
so the Douglas bridge exists.
\end{theorem}

\section{What is earned and what remains}
This note earns a positive candidate for the gamma slot as a renormalized shift-difference feature. It does not prove that the classical gamma term in a chosen explicit-formula convention equals this candidate plus accepted constants and residuals. The remaining analytic task is the exact matching of conventions, constants, pole/completion terms, and tails on the chosen completed response space.

\end{document}
